\subsubsection{Gripper for debris}
For rotating brush it was decided to use DC motor. It allows for the fastest debris collection. \newline
The first part of the module is the mounting to the base and axis. The brush is fixed on the axis.
	\begin{center}
		\begin{tabular}{|p{0.14\linewidth}|p{0.12\linewidth}|p{0.12\linewidth}|}
			{\includegraphics[scale=0.5]{days_L/Gripper/images/01}} & {\includegraphics[scale=0.5]{days_L/Gripper/images/02}} & {\includegraphics[scale=0.5]{days_L/Gripper/images/03}}\\
			\hline & Mount for axis  &
		\end{tabular}
	\end{center}
	Also protection for wiring of the motor was made from a construction of Tetrix plates that were fixed to the base rigidly.
\begin{figure}[H]
		\begin{minipage}[h]{\linewidth}
			\center{\includegraphics[scale=0.5]{days_L/Gripper/images/04}}
			\caption{Protection on DCmotor}
		\end{minipage}
	\end{figure}
	The next element of the module is the brush. The brush should be tough enough so that it doesn't bend during collecting debris. But it should be elastic so when robot hits the wall of field the brush will bend but not broke.\newline
	At first the brush was made of a piece of a plastic bottle. During the test it collect debris well enough but it started to crack.
	\begin{figure}[H]
		\begin{minipage}[h]{1\linewidth}
			\center{\includegraphics[scale=0.5]{days_L/Gripper/images/05}}
			\caption{Brush madde of plastic bottle}
		\end{minipage}
	\end{figure}
	So it was decided to use pieces of rubber hose because it is elastic but more durable than plastic. The pieces of hose were fixed on the aluminium tube from Tetrix set with screws. The tube was instaled on the shaft of the motor instead of the axis. This construction worked well enough. It collected the debris without problems.  
	\begin{figure}[H]
		\begin{minipage}[h]{\linewidth}
			\center{\includegraphics[scale=0.5]{days_L/Gripper/images/06}}
			\caption{Gripper with hose}
		\end{minipage}
	\end{figure}
	Later it was decided to increase speed of rotation and install a transmission with gear ratio 2:1. It was needed for pushing away scoring elements during autonomous period.
   \begin{figure}[H]
		\begin{minipage}[h]{\linewidth}
			\center{\includegraphics[scale=0.5]{days_L/Gripper/images/07}}
			\caption{Transmission with gear}
		\end{minipage}
	\end{figure}   
But this construction didn't worked, because the pieces of the hose were inflexible and got stuck in scoring elements while robot wasn't moving. So we decided to use tubes from a dropper. This also meant we could add more tubes, because they not as wide. This construction worked very well.
   \begin{figure}[H]
		\begin{minipage}[h]{\linewidth}
			\center{\includegraphics[scale=0.5]{days_L/Gripper/images/08}}
			\caption{Gripper with tubes from dropper}
		\end{minipage}
	\end{figure}
	
\fillpage	
